\chapter{The native Ada/SPARK simulation: methods and complete coverage}
\label{sec:seq-sim-spark}

\section{Recorded experiment and its source context}
The native results in this chapter come from
\path{evidence/campaign-20261114-current/runtime/simulation.csv}.
The directory name identifies the campaign taken on 14 November 2026;
it is not a claim that its source hashes match every later checkout.
Its \path{campaign.json}, \path{sources.tar.gz}, native proof logs and
\path{runtime/validation.json} identify the actual source and toolchain.
The archived CSV has SHA-256
\begin{quote}\small\ttfamily
\nolinkurl{d4647eb637d7b2e3520a7bf694b1de8ba2eca7d281e8c41c36d53ba8e1c8a3a8}.
\end{quote}
Tables in this chapter transcribe that record at the stated decimal
precision; they are not newly generated random observations. The Dask
tables elsewhere retain their separate campaign and input hashes.

The native driver evaluates 20 finite parameter cases and samples a
population and a conditional aggregate for each. It then samples six
limiting laws, one population and one conditional law for each
\(r=1,2,3\) at \(c=0.6,a=10/3\). Every batch contains 200,000
scalar observations. Thus there are \(2\cdot20+6=46\) batches and
\(46\cdot200{,}000=9{,}200{,}000\) scalar output observations.
This count includes deterministic zero-variance cases and is not a count
of independent experimental comparisons. The finite grid comprises the
nine requested IQ-150 cases, three \(a=2,k=3\) examples, four cases at
\(c=0\) or \(c=1\), one negative-correlation interior, and three
singular negative endpoints. These counts follow by adding
\(9+3+4+1+3=20\), not from extrapolating the smaller Dask grid.

\section{Scalar aggregate sampling and why its law is correct}
For a population batch the native algorithm sets location \(m=0\)
and scale \(\sigma=\sqrt{d_N/N}\). For a conditional batch it sets
\(m=B_{r,k}\) and
\(\sigma=\sqrt{k(1-c)d_N/(N^2d_r)}\). It draws the scalar
\(m+\sigma H\), rather than constructing all \(N\) scores.
Consequently the work per output is independent of the number of scores
whose average is represented. This shortcut checks the aggregate
sampler and reductions against a supplied law; it does not independently
derive the supplied variance. The recorded matrix and exact-rational
checks described below address that separate calculation.

\begin{proposition}[Scalar simulation law, including finite negative cases]
\label{seq-sim:scalar-law}
Let \(-1/(N-1)\leq c<1\), \(N=r+k\), \(k\geq1\).
An equicorrelated Gaussian battery has population aggregate law
\(N(0,d_N/N)\), and its continuous conditional law at
\(X_0=\cdots=X_{r-1}=a\) is
\[
 N\left(B_{r,k},\frac{k(1-c)d_N}{N^2d_r}\right).
\]
Here \(N(m,0)\) denotes point mass at \(m\). Thus drawing a scalar
with these parameters reproduces the aggregate law even for the finite
negative-correlation cases, without using a negative square root of
\(c\). At \(c=1\), the population law is \(N(0,1)\) and the
conditional law on the equal-score support is point mass at \(a\).
\end{proposition}
\begin{proof}
Proposition~\ref{seqp:domain} constructs the Gaussian vector from its
nonnegative covariance eigenvalues, including the singular endpoint.
The aggregate is a Gaussian linear transformation by
Proposition~\ref{pf:gaussian}, and summing its covariances gives
\(d_N/N\). For the initial block, \(1-c>0\) and
\(d_r=d_N-kc>0\): this is immediate if \(c\geq0\), and follows
from \(d_N\geq0,-kc>0\) if \(c<0\). Hence its covariance
\(R_r\) is invertible and
\(R_r^{-1}\mathbf1_r=\mathbf1_r/d_r\).
The Gaussian regression/residual theorem
\citep[Theorems 3.5.1--3.5.2]{gallager2013}, proved in
Proposition~\ref{pf:conditioning}, gives the future block mean
\((c/d_r)(ra)\mathbf1_k=b_r\mathbf1_k\) and covariance
\[
 (1-c)I_k+\left(c-\frac{rc^2}{d_r}\right)
        \mathbf1_k\mathbf1_k^{\mathsf T}
 =(1-c)I_k+\frac{c(1-c)}{d_r}\mathbf1_k\mathbf1_k^{\mathsf T}.
\]
Its total sum has variance
\(k(1-c)+k^2c(1-c)/d_r=k(1-c)d_N/d_r\).
The fixed initial sum adds no conditional variance; division by \(N^2\)
proves the formula. Its conditional centre is
\((ra+kb_r)/N=B_{r,k}\). This proof uses the full PSD covariance,
not the nonnegative common-factor realization. If \(d_N=0\), the
aggregate has zero variance and zero population mean, hence is zero
almost surely: \(\E[C^2]=0\) implies
\(P(|C|>\epsilon)\leq\E[C^2]/\epsilon^2=0\) for every
\(\epsilon>0\). The displayed conditional centre and variance are
also zero. Finally, when \(c=1\),
\(\E[(X_i-X_j)^2]=0\), so all scores equal one standard-normal
variable almost surely; observing their common value at \(a\) gives
the stated conditional point mass. These are the finite-second-moment
and Gaussian results of \citet{durrett2019,gallager2013}, with each
substitution made explicitly here.
\end{proof}

The limit batches use the already proved limiting means and variances
from \Cref{sec:seq-sim-limits}. They draw those scalar Gaussian laws
directly. Native and Dask limit experiments therefore sample the same
target laws by different implementation paths; neither replaces the
convergence proof with a finite experiment.

\section{Generator, state transitions and moment reductions}
In the archived program, \texttt{Sequential\_Sampling.Draw} reuses
\texttt{Normal\_Support.Normal\_Pair}. The latter uses a 64-bit linear
congruential state transition, projects its high 52 bits to an open-unit
floating-point value, and applies a Box--Muller transformation using the
Ada elementary-function library. Each nondegenerate scalar draw uses
the first output of one pair and discards the second; there is no hidden
cached second draw. The driver resets the state for each batch, increasing
the seed from \texttt{20261115} through \texttt{20261160} in CSV order.
Those are reproducibility identifiers, not calendar dates or proofs of
independence. With zero scale, \texttt{Draw} returns its location and
leaves the state unchanged, so deterministic rows consume no random
normal pairs. The archived source, rather than the later shared-library
wrapper with the same intended algorithm, defines this experiment.

The generator is a finite deterministic mechanism. The normal model is
the reference law against which its outputs are diagnosed; the program's
range and state contracts do not prove exact IID Gaussian randomness.
The native driver stores a count \(M\), sum \(S\) and square sum
\(T\), and reports
\[
 \widehat m=S/M,\qquad
 s^2=\max\left(0,\frac{T-S^2/M}{M-1}\right).
\]
In exact arithmetic,
\(T-S^2/M=\sum_{i=1}^{M}(x_i-\widehat m)^2\geq0\), by expanding
the squares and using \(\sum_i x_i=M\widehat m\). Thus the maximum
does not alter the exact sample variance; its purpose in floating-point
code is to guard cancellation. Both implementations use denominator
\(M-1\). Under the stated ideal IID Gaussian assumption,
Proposition~\ref{sm:normal-sample} proves the unbiasedness, mean SE and
chi-square law for this statistic; it does not certify the rounding error
of the native floating-point accumulator. For example, the deterministic
\(c=1\) rows print a mean about
\(3.33333333333552\), rather than the infinitely precise \(10/3\).
Their zero recorded variance does not remove that finite accumulation
error.

\section{Finite results beyond the requested IQ grid}
The nine native IQ-150 population/conditional pairs are reported in the
next chapter alongside the Dask replication discussion. The table below
supplies all 11 additional native finite pairs, so all 40 finite batches
are covered. Entries are standardized sample means and sample variances,
rounded to six decimal places from the retained CSV; a displayed zero
can mean rounding, whereas the singular rows contain literal CSV zeros.

\begin{table}[htbp]
\centering\small
\setlength{\tabcolsep}{4pt}
\begin{tabular}{rrrrrrrr}
\toprule
\(r\)&\(k\)&\(c\)&\(a\)&\(\widehat m_P\)&\(s_P^2\)&\(\widehat m_C\)&\(s_C^2\)\\
\midrule
1&3&0.6&2&0.001056&0.699618&1.400384&0.208876\\
2&3&0.6&2&-0.001582&0.680839&1.699680&0.101796\\
3&3&0.6&2&-0.000583&0.670954&1.818461&0.060740\\
1&2&0&\(10/3\)&0.000762&0.334988&1.110269&0.221453\\
3&10&0&\(10/3\)&0.000003&0.077044&0.769909&0.059139\\
1&2&1&\(10/3\)&0.001390&0.994377&3.333333&0.000000\\
3&10&1&\(10/3\)&0.002237&1.000936&3.333333&0.000000\\
2&2&-0.25&2&0.000439&0.062565&0.333045&0.052169\\
1&1&-1&2&0.000000&0.000000&0.000000&0.000000\\
1&2&-0.5&2&0.000000&0.000000&0.000000&0.000000\\
3&2&-0.25&2&0.000000&0.000000&0.000000&0.000000\\
\bottomrule
\end{tabular}
\caption{Additional native finite cases: 200,000 observations in each
population (P) or conditional (C) batch. Source: the native CSV identified
in \Cref{sec:seq-sim-spark}. These cases are outside the Dask IQ-150 grid.}
\label{tab:seq-native-extra}
\end{table}

The first three rows test a different starting level, without changing
the population law. Substituting \(a=2,c=3/5,k=3\) in the proved
formulas gives conditional centres \(7/5,17/10,20/11\), population
variances \(7/10,17/25,2/3\), and conditional variances
\(21/100,51/500,2/33\). These targets explain both the shift in
conditional means and the decreasing conditional variances across this
fixed-\(k\), increasing-\(r\) comparison. For example,
\(d_6=4,d_3=11/5\) gives
\(B_{3,3}=3\cdot4\cdot2/(6\cdot11/5)=20/11\) and
\(v_C=3(2/5)4/[36(11/5)]=2/33\). The other entries follow by
the same explicit substitution, with \((d_N,d_r)=(14/5,1)\) and
\((17/5,8/5)\), respectively.

At \(c=0\), direct substitution gives
\(B_{r,k}=ra/N,v_P=1/N,v_C=k/N^2\): later scores have zero
prediction, but the initial scores remain included. The two table rows
therefore have conditional target means \(10/9\) and \(10/13\),
and conditional target variances \(2/9\) and \(10/169\).
At \(c=1\), conditioning reveals every future score, producing zero
conditional spread while whole random batteries retain unit population
variance, as proved above. The negative interior row has
\(N=4,d_N=1/4,d_r=3/4\), so
\(B=1/3,v_P=1/16,v_C=5/96\), consistent with the recorded
\(0.333045,0.062565,0.052169\). The last three rows have
\(d_N=0\); their exact zero aggregates are consequences of the
covariance singularity, not evidence of a malfunctioning random sampler.
Fixed negative \(c\) has no admissible infinite-size continuation,
by Proposition~\ref{seqp:domain}; no such limit is reported.

\section{Native limits and a direct comparison with Dask}
\begin{table}[htbp]
\centering\small
\setlength{\tabcolsep}{4pt}
\begin{tabular}{rlrrrrrr}
\toprule
\(r\)&Law&\(m_*\)&\(\widehat m_S\)&\(\widehat m_D\)&\(v_*\)&\(s_S^2\)&\(s_D^2\)\\
\midrule
1&Pop.&0.000000&-0.000834&-0.000987&0.600000&0.594883&0.599413\\
1&Cond.&2.000000&2.001589&1.999376&0.240000&0.241082&0.239765\\
2&Pop.&0.000000&-0.000792&0.000076&0.600000&0.600198&0.599842\\
2&Cond.&2.500000&2.500597&2.500038&0.150000&0.150202&0.149961\\
3&Pop.&0.000000&0.001884&0.000400&0.600000&0.600157&0.601390\\
3&Cond.&2.727273&2.726864&2.727443&0.109091&0.109373&0.109344\\
\bottomrule
\end{tabular}
\caption{All six native limiting batches (S), compared with the pooled
Dask batches (D). Centres and variances are standardized; stars denote
analytic targets. Native \(M=200{,}000\); Dask \(M=4{,}000{,}000\).
The Dask source is the retained moments CSV.}
\label{tab:seq-native-limits}
\end{table}

The limiting conditional targets are exactly
\((m_*,v_*)=(2,6/25),(5/2,3/20),(30/11,6/55)\), from the
limit formulas already proved. The population targets are \((0,3/5)\).
Thus the six table rows concern specified laws, not an estimated
infinite number of tests. The native \(r=1\) population variance
\(0.594883\) is visibly farther from \(0.6\) than its pooled Dask
counterpart \(0.599413\), whereas a larger Dask sample need not be
closer for every particular seed: the \(r=3\) population variance is
one example in the opposite direction. Sample size controls the reference
sampling distribution, not the deterministic ordering of two realized
errors.

For a finite comparison, with \(r=1,k=10\), the native conditional
IQ centre and SD are approximately \(131.819349\) and \(7.193671\),
against Dask's \(131.809593\) and \(7.213895\); the respective
analytic targets are \(1450/11\) and \(\sqrt{6300/121}\).
Both algorithms therefore report outcome spreads near 7.2 IQ points;
the much smaller uncertainty in a simulated centre concerns a different
quantity. The ideal mean-estimation SE is \(\sqrt{v/M}\). For a
fixed positive-variance law the native-to-Dask SE ratio is therefore
\(\sqrt{4{,}000{,}000/200{,}000}=\sqrt{20}\), directly from
Proposition~\ref{sm:normal-sample}.

No cross-implementation independence is inferred from their different
generators. If \(\widehat m_S,\widehat m_D\) estimate the same centre,
covariance bilinearity gives
\[
 \Var(\widehat m_S-\widehat m_D)
 =\Var(\widehat m_S)+\Var(\widehat m_D)
  -2\Cov(\widehat m_S,\widehat m_D).
\]
Only under an additional ideal independence assumption does this become
\(v(1/200{,}000+1/4{,}000{,}000)\). The displayed comparisons are
descriptive recorded differences; they do not silently use that
unestablished cross-generator assumption to assign a paired confidence
level.

\section{What was checked natively and what remains trusted}
The recorded validator checks all 20 calculation cases against exact
rational moment formulas and full covariance-matrix calculations,
including a pseudoinverse comparison at singular boundaries. It checks
all 46 simulation rows against their own population or conditional
targets, treating zero-variance rows separately. Its recorded maximum
absolute mean discrepancy is \(1.525037766488471\) theoretical SEs
over the positive-variance rows. It also reports 18 identity and finite-PMF
checks, whose largest absolute residual is
\(1.77635683940025\times10^{-15}\). These are finite numerical
diagnostics, with tolerances in the source and CSV; no certified
transcendental or floating-point error bound follows from the small
observed residual.

For positive variances, the validator uses a two-sided normal mean
diagnostic and a two-sided chi-square variance diagnostic, each with
ideal-model rejection probability \(10^{-10}\), apart from its explicit
numerical tolerances. Their reference laws are proved in
Proposition~\ref{sm:normal-sample}. Applying the union bound from
Proposition~\ref{sm:multiple} to at most \(2\cdot46\) diagnostics
gives an ideal-law false-alarm upper bound \(92\cdot10^{-10}\),
without requiring independent rows. Zero-variance rows instead use
absolute tolerances. This deliberately broad bound concerns the stated
ideal sampling laws; it is not a proof that a finite-state PRNG generates
IID normal observations.

The source-pinned campaign records 430 discharged checks (124 flow,
306 prover), with zero unproved or justified. Its 11 sequential interfaces
have 139 clauses; all meet the five-clause floor and eight meet ten.
Declaration counts do not establish contract independence or adequacy.
The checked square-root and elementary-function/I/O boundaries remain
trusted; the CLI is outside SPARK. Subsequent shared-library refactoring
changed source hashes, so these historical counts do not verify current
Ada. Dask/Deal tests, lint, symbolic checks, written proofs and recorded
native proofs each retain their stated scope. None establishes Gaussian
PRNG distributional validity or empirical IQ-model validity.
